\documentclass{article}%
\usepackage[T1]{fontenc}%
\usepackage[utf8]{inputenc}%
\usepackage{lmodern}%
\usepackage{textcomp}%
\usepackage{lastpage}%
\usepackage{geometry}%
\geometry{margin=2.2cm}%
%
\usepackage{amsmath}%
\usepackage{amssymb}%
\usepackage{booktabs}%
\usepackage{longtable}%
\usepackage{array}%
\usepackage{hyperref}%
\hypersetup{hidelinks}%
\title{ITA 2027 {-} Química (claude\_sonnet\_5\_5) x Poliedro\\[4pt]\large Comparação com o gabarito Poliedro Resolve (ITA 2027, 2ª fase, dia 2 {-} Química)\thanks{Gerado com gpt-resolve: \url{https://github.com/lgabs/gpt-resolve}}}%
\date{09/11/2026}%
%
\begin{document}%
\normalsize%
\maketitle%
\section*{Resumo}\textbf{Nota média: 9,8/10} (33 de 34 itens com a resolução correta; cada item vale o mesmo dentro da questão).\par\medskip%
\begin{longtable}{@{}c c p{0.78\textwidth}@{}}\toprule Questão & Nota & Resumo \\ \midrule
1 & 10,0 & A resolução da questão 1 coincide com o gabarito do Poliedro nos quatro itens: (a) C5H12 + 8 O2 $\to$ 5 CO2 + 6 H2O; (b) 1,6 mol de O2 e 6,4 mol de N2; (c) T = 2300,84 K; (d) redução de $\approx$ 37\% no combustível, com 0,1267 mol de pentano em mistura pobre. Refiz todas as contas em Python e elas batem. Não há erro na resolução nem no gabarito; o que ficou de dúvida são só hipóteses implícitas que os dois adotam. \\
2 & 10,0 & A resolução está correta e coincide com o gabarito do Poliedro. Os dois chegam a $\pi$ = 5,008/(10 - x) atm entre 0 e 6 dm, subindo de 0,50 a 1,25 atm, e a um patamar de 1,252 atm entre 6 e 9 dm, quando a solução atinge a saturação de 0,05 mol/L. Recalculei todas as contas e os pontos tabelados conferem. As hipóteses adotadas são as mesmas: soluto não eletrólito e solução ideal. Não há divergência. \\
3 & 10,0 & Os quatro itens da resolução estão certos e coincidem com o gabarito do Poliedro, que também está correto. As respostas são: X = 2{-}metilbut{-}2{-}eno (Zaitsev); Z = 2{-}bromo{-}2{-}metilbutano (Markovnikov); G = 2{-}bromo{-}3{-}metilbutano (radicalar, anti{-}Markovnikov); J = 2{-}metilbutano{-}2,3{-}diol (di{-}hidroxilação sin com KMnO4 diluído, a frio). Os mecanismos, as fórmulas moleculares e as observações de estereoquímica (G e J racêmicos) também estão corretos. A única ressalva vem do enunciado: ele usa H2O2 como 'peróxido' para o efeito Kharasch, o que o próprio gabarito comenta. A resolução segue a mesma leitura, então a ressalva não afeta a nota. \\
4 & 10,0 & Conferi o enunciado transcrito com o original (página 6 do PDF) e ele está correto. A resolução bate com o gabarito do Poliedro nos quatro itens: a) 40 g; b) 75\%; c) 180 g/mol; d) 3 mol/L. Refiz as contas e todas estão certas: n\_cát = 1 mol, dividido em 0,6, 0,3 e 0,1 mol, e n\_X$^{-}$ = 1,5 mol. Não há nenhuma divergência. A resolução está correta e completa. \\
5 & 10,0 & Nos itens a (N2 + 8H+ + 6e- $\to$ 2NH4+) e b (Nox +2, +5, 0, -3, +3 e +1 médio para o N2O), a resolução coincide com o gabarito e está correta. No item c, a ordem das espécies é a mesma, N2 $\to$ NH4+ $\to$ NO2- $\to$ NO3- $\to$ NO2- $\to$ NO $\to$ N2O. A diferença é que o gabarito fecha o ciclo com $\to$ N2, e a resolução para em N2O porque o enunciado fala em desnitrificação 'parcial'. Ela só apresenta o fechamento em N2 como leitura alternativa. Considero as duas sequências aceitáveis, e todas as semirreações conferem em átomos e carga. O único risco de perder nota é o item c, se o corretor exigir o N2 no fim da resposta principal. \\
6 & 10,0 & A resolução está correta e completa. A reação 1 (Li + H2O) e a reação 3 (reforma a vapor) coincidem com o gabarito em estados, estequiometria e condições. Na reação 2, a resolução usa a pirólise do CH4 no lugar da eletrólise da água do gabarito. O enunciado permite reagentes sozinhos e o gabarito só dá exemplos, então a escolha é válida, e a termodinâmica citada confere ($\Delta$S $\approx$ +81 J/K, T* $\approx$ 925 K, $\Delta$H $\approx$ +75 kJ/mol). Por causa dessa escolha, a resposta do item iv é diferente: a resolução aponta a pirólise (2 mol H2/mol de reagente), e o gabarito, a reforma (1,5). Cada uma está certa para o próprio conjunto de reações, usando a mesma convenção de contagem. A da resolução é inclusive o maior valor entre as reações possíveis. Ela ainda discute a base por mol de CH4, em que a reforma vence. Não encontrei erro nem na resolução nem no gabarito, só a imprecisão de 'alta pressão' na reforma, do lado do gabarito. \\
7 & 10,0 & A resolução coincide com o gabarito do Poliedro nos quatro itens. No item a, as equações das duas etapas e as condições (oxidante forte acidificado, a quente e sob refluxo; esterificação de Fischer com H2SO4 e aquecimento) são as mesmas. Os itens b ($\Delta$H$^\circ$ = -496,2 kJ/mol), c ($\Delta$S$^\circ$ = -127,2 J/(mol$\cdot$K)) e d ($\Delta$G$^\circ$ $\approx$ -458,3 kJ/mol) foram recalculados em Python e estão corretos. O enunciado foi conferido no PDF original (página 9 impressa). Não há erros; só observações menores de redação e notação. \\
8 & 10,0 & A resolução e o gabarito coincidem nos três itens, e os dois estão corretos. O enunciado conferido no PDF (pág. 10 da prova) bate com a transcrição. a) S e Se têm 2 elétrons desemparelhados cada (p$^{4}$); K tem 1 elétron de valência (4s$^{1}$). b) H2S tem 2 pares não ligantes no S; SeF4 tem 1 par não ligante no Se. c) H2S é angular e SeF4 é gangorra. A justificativa pelo modelo VSEPR está certa, inclusive a posição equatorial do par livre no SeF4. Não há erro que mude a nota. \\
9 & 10,0 & A resolução está correta e bate com o gabarito do Poliedro nos dois itens: a) {[}HX{]} = 8$\cdot$10$^{-}$$^3$ mol/L; b) {[}HY{]} = 1,2$\cdot$10$^{-}$$^2$ mol/L. O caminho é o mesmo do gabarito: soma de mols pela base gasta, hidrólise de Y$^{-}$ no equivalente com pH 9 dando {[}Y$^{-}${]} = 10$^{-}$$^2$ mol/L, e correção da diluição de 120 mL para 100 mL. Conferi as contas sem aproximações em python (pH no equivalente $\approx$ 8,9998). O enunciado transcrito confere com o PDF original. A única falha é um erro de digitação no gabarito ({[}HX{]} no lugar de {[}HY{]}), que não afeta a resposta. \\
10 & 8,0 & Das cinco partes, a resolução concorda com o gabarito em quatro: a(1), a(3), a(4) e b (proteção catódica nos quatro casos). Diverge em a(2), onde aponta a região alcalina X como ânodo e a ácida Y como cátodo, com base no Nernst para Fe(OH)2. As contas estão certas (E\_Fe $\approx$ -0,1 - 0,059$\cdot$pH), mas a interpretação ignora a passivação do Fe em pH alto e a dissolução ativa em meio ácido. Pela leitura padrão de corrosão, o gabarito está certo: X (pH alto) é a área catódica e Y (pH baixo) a anódica. A própria resolução reconhece essa inversão na ressalva, mas não a adota. Resultado: perde o subitem a(2) e acerta o restante. \\
\bottomrule\end{longtable}%
\section*{Principais pontos de dúvida}%
\begin{itemize}%
\item%
\textbf{Q5} \textit{(solution)} No item c, a resposta em destaque da resolução termina em N2O, e o gabarito termina em N2 para fechar o ciclo. A resolução apresenta a sequência com $\to$ N2 só como 'leitura alternativa'. Um corretor que espere o ciclo fechado pode tirar parte da nota pela falta do N2 na resposta principal ou por ela dar duas respostas. As duas sequências são quimicamente corretas, então o risco é pequeno, mas um professor deve decidir.%
\item%
\textbf{Q10} \textit{(solution)} No item a(2) a resolução trocou cátodo e ânodo. Ela usou só o potencial de equilíbrio de Nernst com Fe(OH)2 e deixou de fora a passivação do Fe em meio alcalino e a dissolução ativa em meio ácido. Na leitura padrão de corrosão, a região de pH baixo (Y) é o ânodo, como dá o gabarito. A resolução cita essa inversão na ressalva, mas fica com a resposta contrária, o que deve custar a nota do subitem.%
\end{itemize}%
\section*{Questão 1}%
\href{https://poliedroresolve.sistemapoliedro.com.br/vestibulares/ita/2027/ita-2-fase-dia-2-21-10-2026-2-fase-dia-2/dissertativa-1-quimica-2-fase-dia-2-ita-2027}{Resolução do gabarito}\par\smallskip%
\begin{longtable}{@{}p{0.06\textwidth} p{0.33\textwidth} p{0.33\textwidth} p{0.1\textwidth} p{0.08\textwidth}@{}}\toprule Item & Resolução & Gabarito & Veredito & Certo \\ \midrule
a & C5H12 + 8 O2 $\to$ 5 CO2 + 6 H2O ($\Delta$Hc = -3440 kJ/mol) & C5H12 + 8 O2 $\to$ 5 CO2 + 6 H2O (com ar: C5H12 + 8 O2 + 32 N2 $\to$ 5 CO2 + 6 H2O + 32 N2) & igual & ambos \\
b & n\_ar = 8,0 mol; n(O2) = 1,6 mol; n(N2) = 6,4 mol & n\_ar = 8 mol; O2 = 1,6 mol; N2 = 6,4 mol & igual & ambos \\
c & 0,20 mol de pentano; Q = 688 kJ; 8,6 mol de gases; $\Delta$T = 2000 K; T = 2300,84 K & 0,2 mol de pentano; $\Delta$T = 2000 K $\Rightarrow$ T\_f = 2300,84 K & igual & ambos \\
d & Mistura pobre: x $\approx$ 0,127 mol de pentano ($\approx$ 63\% do estequiométrico), ou seja, reduzir o combustível em $\approx$ 37\% & Reduzir o combustível em $\approx$ 37\% (0,1267 mol de pentano, mistura pobre) & igual & ambos \\
\bottomrule\end{longtable}%
\begin{itemize}%
\item%
\textbf{a)} Mesma equação balanceada (C: 5; H: 12; O: 16). O gabarito também escreve a versão com o N2 do ar, mas isso é só um complemento e não muda a resposta.%
\item%
\textbf{b)} Conferido: n = PV/RT = (10\^{}5 $\cdot$ 0,200)/(8,31 $\cdot$ 300,84) = 8,00006 $\approx$ 8,0 mol. Com 20\%/80\%, dá 1,6 mol de O2 e 6,4 mol de N2, com uma casa decimal, como o enunciado pede.%
\item%
\textbf{c)} Conferido: 0,2 $\cdot$ 3440$\cdot$10\^{}3 / (8,6 $\cdot$ 40) = 2000 K, logo T = 2300,84 K. Os gases finais são 6,4 N2 + 1,0 CO2 + 1,2 H2O = 8,6 mol, exatamente como no gabarito.%
\item%
\textbf{d)} Conferido: 3440000x = 40(8 + 3x)$\cdot$1300 dá x = 416000/3284000 = 0,12667 mol, e (0,2 - 0,12667)/0,2 = 36,66\% $\approx$ 37\% de redução. A resolução e o gabarito usam o mesmo modelo: o O2 em excesso entra na mistura e n\_f = 8 + 3x. A resolução também justifica por que descartou a mistura rica: refiz a conta e ela daria $\approx$ 4,8 mol de pentano no total, dos quais $\approx$ 4,6 mol não queimam, o que bate com o que a resolução diz.%
\end{itemize}%
\paragraph{Pontos de dúvida}%
\begin{itemize}%
\item%
\textit{(statement)} Para fechar a conta da mistura rica, seria preciso supor que o pentano não queimado tem o mesmo Cp de 40 J/(K$\cdot$mol). O enunciado só cita CO2, H2O, N2 e O2. Isso não pesa, porque a 'economia de combustível' e a 'combustão completa' já excluem a mistura rica.%
\item%
\textit{(solution)} Hipóteses implícitas, todas razoáveis e iguais às do gabarito: a água fica como vapor e entra na contagem de mols dos gases; o Cp é o mesmo para qualquer composição; só os produtos absorvem o calor; e o O2 em excesso entra no n\_f do item d.%
\item%
\textit{(solution)} Para o item d, a resolução dá ao mesmo tempo '$\approx$ 63\% do estequiométrico' e 'redução de $\approx$ 37\%'. A resposta pedida, que é a alteração em porcentagem inteira, é a redução de 37\%, e ela está declarada claramente.%
\end{itemize}%
\section*{Questão 2}%
\href{https://poliedroresolve.sistemapoliedro.com.br/vestibulares/ita/2027/ita-2-fase-dia-2-21-10-2026-2-fase-dia-2/dissertativa-2-quimica-2-fase-dia-2-ita-2027}{Resolução do gabarito}\par\smallskip%
\begin{longtable}{@{}p{0.06\textwidth} p{0.33\textwidth} p{0.33\textwidth} p{0.1\textwidth} p{0.08\textwidth}@{}}\toprule Item & Resolução & Gabarito & Veredito & Certo \\ \midrule
Questão 2 (gráfico $\pi$ $\times$ deslocamento, 0 a 9 dm) & $\pi$(x) = 0,2$\cdot$0,0821$\cdot$305/(10 - x) = 5,008/(10 - x) atm para 0 $\leq$ x $\leq$ 6 dm (0,50 $\to$ 1,25 atm, hipérbole). Em x = 6 dm a solução satura (0,05 mol/L) e $\pi$ fica constante em $\approx$ 1,25 atm até x = 9 dm. Tabela: 0,50; 0,56; 0,63; 0,72; 0,83; 1,00; 1,25; 1,25; 1,25; 1,25. & $\pi$(x) = 5,0081/(10 - x) atm para 0 $\leq$ x $\leq$ 6 dm ($\pi$(0) $\approx$ 0,50 atm até $\approx$ 1,25 atm); de 6 a 9 dm, $\pi$ = 1,252 atm constante (saturação a 0,05 mol/L). Gráfico: hipérbole crescente de 0 a 6 e patamar de 6 a 9. & igual & ambos \\
\bottomrule\end{longtable}%
\begin{itemize}%
\item%
\textbf{Questão 2 (gráfico $\pi$ $\times$ deslocamento, 0 a 9 dm))} Refiz as contas: RT = 0,0821$\cdot$305 = 25,04 atm$\cdot$L/mol; nRT = 5,0081 atm$\cdot$L; V = 10 - x (L), porque a área é 1 dm$^2$. A saturação ocorre em 0,2/(10 - x) = 4/80 = 0,05 mol/L, ou seja, x = 6 dm, com $\pi$\_sat = 0,05$\cdot$25,04 = 1,252 atm. Os valores da tabela batem com os do gabarito ponto a ponto: 0,501; 0,556; 0,626; 0,715 ($\approx$ 0,72); 0,835 ($\approx$ 0,83); 1,002; 1,252. A resolução também confere que c$_{0}$ = 1,6 g/L é menor que 4 g/L, adota as mesmas hipóteses do gabarito (i = 1, solução ideal) e traça o gráfico com a mesma forma: hipérbole de 0 a 6 dm e patamar de 6 a 9 dm.%
\end{itemize}%
\paragraph{Pontos de dúvida}%
\begin{itemize}%
\item%
\textit{(statement)} O enunciado não diz se o soluto é eletrólito. A resolução e o gabarito admitem i = 1 (não eletrólito). Se o soluto fosse eletrólito, todos os valores de $\pi$ seriam multiplicados por i, mas a forma do gráfico (hipérbole + patamar em x = 6 dm) não mudaria. É uma hipótese razoável, e as duas versões a declaram.%
\item%
\textit{(statement)} O enunciado chama de $\pi$ tanto a sobrepressão aplicada ao pistão quanto a pressão osmótica. Ambas as resoluções identificam as duas porque o processo é reversível (quase{-}estático). Também desprezam o volume do precipitado, a diferença de coluna hidrostática entre os compartimentos e a possibilidade de o êmbolo, em x $\approx$ 9 dm, chegar perto da altura da membrana mostrada na figura. São simplificações implícitas e aceitáveis.%
\item%
\textit{(key)} O gráfico do gabarito é uma imagem externa à qual não tive acesso. Conferi apenas a função por partes e os valores tabelados, que coincidem com os da resolução.%
\end{itemize}%
\section*{Questão 3}%
\href{https://poliedroresolve.sistemapoliedro.com.br/vestibulares/ita/2027/ita-2-fase-dia-2-21-10-2026-2-fase-dia-2/dissertativa-3-quimica-2-fase-dia-2-ita-2027}{Resolução do gabarito}\par\smallskip%
\begin{longtable}{@{}p{0.06\textwidth} p{0.33\textwidth} p{0.33\textwidth} p{0.1\textwidth} p{0.08\textwidth}@{}}\toprule Item & Resolução & Gabarito & Veredito & Certo \\ \midrule
a & X = (CH3)2C=CH-CH3 (2{-}metilbut{-}2{-}eno), produto de Zaitsev da desidratação E1 & X = 2{-}metilbut{-}2{-}eno, (CH3)2C=CH-CH3 & igual & ambos \\
b & Z = (CH3)2CBr-CH2CH3 (2{-}bromo{-}2{-}metilbutano), adição de Markovnikov & Z = 2{-}bromo{-}2{-}metilbutano, (CH3)2CBr-CH2CH3 & igual & ambos \\
c & G = (CH3)2CH-CHBr-CH3 (2{-}bromo{-}3{-}metilbutano), adição radicalar anti{-}Markovnikov, mistura racêmica & G = 2{-}bromo{-}3{-}metilbutano, (CH3)2CH-CHBr-CH3 (anti{-}Markovnikov) & igual & ambos \\
d & J = (CH3)2C(OH)-CH(OH)-CH3 (2{-}metilbutano{-}2,3{-}diol), di{-}hidroxilação sin de Baeyer & J = 2{-}metilbutano{-}2,3{-}diol, (CH3)2C(OH)-CH(OH)-CH3 & igual & ambos \\
\bottomrule\end{longtable}%
\begin{itemize}%
\item%
\textbf{a)} Resolução e gabarito dão a mesma resposta. Na desidratação do álcool terciário 2{-}metilbutan{-}2{-}ol com H2SO4/$\Delta$, o carbocátion terciário perde um H. Pela regra de Zaitsev predomina o alceno trissubstituído, o 2{-}metilbut{-}2{-}eno. A resolução cita corretamente o 2{-}metilbut{-}1{-}eno como produto minoritário.%
\item%
\textbf{b)} Resolução e gabarito dão a mesma resposta. É uma adição eletrofílica iônica: o H+ entra no C3 e forma o carbocátion terciário em C2, que recebe o Br-. Fórmula C5H11Br, correta.%
\item%
\textbf{c)} Resolução e gabarito dão a mesma resposta. O Br? ataca o C3 e gera o radical terciário em C2, que depois abstrai H de HBr. A observação de que o C do CHBr é estereocentro e sai como racemato está correta. Assim como o gabarito, a resolução segue a intenção da banca de que o 'peróxido (H2O2)' induz o efeito Kharasch.%
\item%
\textbf{d)} Resolução e gabarito dão a mesma resposta. Com KMnO4 diluído, em meio básico e a frio, ocorre di{-}hidroxilação sin sem clivagem da cadeia. A análise estereoquímica também está correta: só o C3 é estereocentro, pois o C2 tem dois CH3, e J sai como par racêmico, sem diastereoisômeros. Fórmula C5H12O2, correta.%
\end{itemize}%
\paragraph{Pontos de dúvida}%
\begin{itemize}%
\item%
\textit{(statement)} O enunciado indica 'Peróxido (H2O2)' para a adição radicalar de HBr. O efeito Kharasch costuma ser feito com peróxidos orgânicos (ROOR), e em H2O2 aquoso a via iônica (Markovnikov) poderia competir. O próprio gabarito do Poliedro registra essa imprecisão e entende que a banca queria o produto anti{-}Markovnikov. A resolução adota a mesma leitura (escreve RO--OR no mecanismo), então a nota não é afetada.%
\item%
\textit{(solution)} A resolução trata X apenas como o produto majoritário (Zaitsev) e não discute a mistura com o 2{-}metilbut{-}1{-}eno nos itens seguintes. É a convenção esperada pela banca e não muda a nota.%
\end{itemize}%
\section*{Questão 4}%
\href{https://poliedroresolve.sistemapoliedro.com.br/vestibulares/ita/2027/ita-2-fase-dia-2-21-10-2026-2-fase-dia-2/dissertativa-4-quimica-2-fase-dia-2-ita-2027}{Resolução do gabarito}\par\smallskip%
\begin{longtable}{@{}p{0.06\textwidth} p{0.33\textwidth} p{0.33\textwidth} p{0.1\textwidth} p{0.08\textwidth}@{}}\toprule Item & Resolução & Gabarito & Veredito & Certo \\ \midrule
a & m = 40 g & m = 40 g & igual & ambos \\
b & P = 75\% & P = 75\% & igual & ambos \\
c & MM = 180 g/mol & MM = 180 g/mol & igual & ambos \\
d & {[}X$^{-}${]} = 3 mol/L & {[}X$^{-}${]} = 3 mol/L & igual & ambos \\
\bottomrule\end{longtable}%
\begin{itemize}%
\item%
\textbf{a)} n\_cátions = 2 $\times$ 0,5 = 1 mol. Na proporção 6:3:1, n\_Q = 0,6 mol. Massa de QX puro = 0,6 $\times$ 60 = 36 g. Com 90\% de pureza, m = 36/0,9 = 40 g. Resolução e gabarito fazem a mesma conta.%
\item%
\textbf{b)} n\_R = 0,3 mol, então a massa de RX2 puro é 0,3 $\times$ 150 = 45 g. P = 45/60 = 0,75, ou seja, 75\%.%
\item%
\textbf{c)} Massa de TX3 puro = 0,9 $\times$ 20 = 18 g, que corresponde a n\_T = 0,1 mol. MM = 18/0,1 = 180 g/mol.%
\item%
\textbf{d)} n\_X$^{-}$ = 0,6 + 2(0,3) + 3(0,1) = 1,5 mol, e 1,5/0,5 = 3 mol/L. A eletroneutralidade confere: a carga positiva também soma 1,5 mol.%
\end{itemize}%
\paragraph{Pontos de dúvida}%
\begin{itemize}%
\item%
\textit{(statement)} O enunciado não diz que as impurezas são inertes, ou seja, que não fornecem Q?, R$^2$?, T$^3$? nem X$^{-}$. É uma hipótese implícita padrão. A resolução a declara de forma explícita e o gabarito a usa sem comentar.%
\item%
\textit{(statement)} 'Concentração total de cátions = 2 mol/L' foi lida como a soma das concentrações em quantidade de matéria (mol), não em equivalentes de carga. Resolução e gabarito adotam a mesma leitura, que é a natural.%
\end{itemize}%
\section*{Questão 5}%
\href{https://poliedroresolve.sistemapoliedro.com.br/vestibulares/ita/2027/ita-2-fase-dia-2-21-10-2026-2-fase-dia-2/dissertativa-5-quimica-2-fase-dia-2-ita-2027}{Resolução do gabarito}\par\smallskip%
\begin{longtable}{@{}p{0.06\textwidth} p{0.33\textwidth} p{0.33\textwidth} p{0.1\textwidth} p{0.08\textwidth}@{}}\toprule Item & Resolução & Gabarito & Veredito & Certo \\ \midrule
a & N2 + 8 H+ + 6 e- $\to$ 2 NH4+ & N2 + 8 H+ + 6 e- $\to$ 2 NH4+ & igual & ambos \\
b & NO +2; NO3- +5; N2 0; NH4+ -3; NO2- +3; N2O +1 (médio) & NO +2; NO3- +5; N2 0; NH4+ -3; NO2- +3; N2O +1 (médio) & igual & ambos \\
c & N2 $\to$ NH4+ $\to$ NO2- $\to$ NO3- $\to$ NO2- $\to$ NO $\to$ N2O. A resolução trata a desnitrificação como parcial. Numa observação, diz que, se ela for completa, o ciclo fecha com $\to$ N2. & N2 $\to$ NH4+ $\to$ NO2- $\to$ NO3- $\to$ NO2- $\to$ NO $\to$ N2O $\to$ N2 (o N2 do ar fecha o ciclo) & diferente & ambos \\
\bottomrule\end{longtable}%
\begin{itemize}%
\item%
\textbf{a)} São idênticas. Conferi átomos e carga: há 2 N e 8 H em cada lado, e a carga é +2 nos dois lados. Cada N passa de 0 a -3, o que dá os 6 e-.%
\item%
\textbf{b)} Os valores são os mesmos. A explicação extra da resolução sobre o N2O também está certa. Na estrutura N$\equiv$N+--O-, o N terminal tem Nox 0 e o central +2. Na estrutura -N=N+=O, eles valem -1 e +3. Nas duas a soma dos N é +2, e por isso só o Nox médio +1 é bem definido.%
\item%
\textbf{c)} A única diferença é o último passo, N2O $\to$ N2. O gabarito o inclui para fechar o ciclo. A resposta principal da resolução para em N2O, porque a introdução do enunciado fala em 'desnitrificação parcial'. Mesmo assim, a resolução escreve explicitamente a sequência do gabarito como leitura alternativa. Tirando isso, as sequências são iguais: fixação N2$\to$NH4+, nitrificação NH4+$\to$NO2-$\to$NO3- e desnitrificação NO3-$\to$NO2-$\to$NO$\to$N2O. Os Nox seguem 0 $\to$ -3 $\to$ +3 $\to$ +5 $\to$ +3 $\to$ +2 $\to$ +1, como deveriam. Conferi todas as semirreações da resolução, inclusive N2O + 2H+ + 2e- $\to$ N2 + H2O, e estão balanceadas em átomos e carga. As duas sequências usam as seis espécies e são 'compatíveis com o ciclo do nitrogênio'. A do gabarito é a leitura mais natural para um ciclo fechado, e a da resolução tem base na palavra 'parcial' do enunciado. Por isso considero as duas aceitáveis.%
\end{itemize}%
\paragraph{Pontos de dúvida}%
\begin{itemize}%
\item%
\textit{(solution)} No item c, a resposta em destaque da resolução termina em N2O, e o gabarito termina em N2 para fechar o ciclo. A resolução apresenta a sequência com $\to$ N2 só como 'leitura alternativa'. Um corretor que espere o ciclo fechado pode tirar parte da nota pela falta do N2 na resposta principal ou por ela dar duas respostas. As duas sequências são quimicamente corretas, então o risco é pequeno, mas um professor deve decidir.%
\item%
\textit{(statement)} O enunciado é ambíguo. A introdução fala em 'desnitrificação parcial', o que permite parar em N2O. Já o item c pede uma sequência 'compatível com o ciclo do nitrogênio' e cita apenas 'desnitrificação', o que sugere voltar ao N2. Além disso, a frase 'uma mesma espécie intermediária poderá aparecer mais de uma vez' só cobre o NO2-. Repetir o N2 no fim não é uma intermediária repetida, e sim o fechamento do ciclo.%
\item%
\textit{(key)} A resolução transcrita do gabarito tem um '5' solto no final, que parece ser o número da questão ou um resto da diagramação. Isso não afeta o conteúdo.%
\end{itemize}%
\section*{Questão 6}%
\href{https://poliedroresolve.sistemapoliedro.com.br/vestibulares/ita/2027/ita-2-fase-dia-2-21-10-2026-2-fase-dia-2/dissertativa-6-quimica-2-fase-dia-2-ita-2027}{Resolução do gabarito}\par\smallskip%
\begin{longtable}{@{}p{0.06\textwidth} p{0.33\textwidth} p{0.33\textwidth} p{0.1\textwidth} p{0.08\textwidth}@{}}\toprule Item & Resolução & Gabarito & Veredito & Certo \\ \midrule
Reação 1: Li + H2O (estados, estequiometria, condições) & 2 Li(s) + 2 H2O(?) $\to$ 2 LiOH(aq) + H2(g); ocorre à temperatura e pressão ambiente, sem catalisador, adicionando o Li à água em excesso (reação exotérmica) & Li(s) + H2O(l) $\to$ LiOH(aq) + 1/2 H2(g) (o mesmo que 2Li + 2H2O $\to$ 2LiOH + H2); basta o contato direto do Li com a água & equivalente & ambos \\
Reação 2 (escolha livre) & Pirólise do metano: CH4(g) $\to$ C(s) + 2 H2(g), $\Delta$H$^\circ$ $\approx$ +75 kJ/mol; ocorre a \textasciitilde{}800--1200 $^\circ$C, sem O2, com pressão baixa (Ni/Fe opcional como catalisador) & Eletrólise da água: H2O(l) $\to$ H2(g) + 1/2 O2(g); precisa de um eletrólito (p. ex. LiOH), eletrodos inertes e corrente contínua & diferente & ambos \\
Reação 3: reforma a vapor do CH4 & CH4(g) + H2O(g) $\to$ CO(g) + 3 H2(g), $\Delta$H$^\circ$ $\approx$ +206 kJ/mol; 700--1000 $^\circ$C, catalisador de Ni/Al2O3, pressão não muito alta (3--25 bar na indústria), vapor em excesso & CH4(g) + H2O(g) $\to$ CO(g) + 3 H2(g); alta temperatura, alta pressão e catalisador adequado & igual & ambos \\
iv) Reação que rende mais H2 por mol de reagente (reação completa) & Pirólise do CH4: 2 mol H2/mol de reagente (reforma 1,5; Li + H2O 0,25). A resolução observa ainda que, contando só o CH4, a reforma rende mais (3 contra 2) & Reforma a vapor do CH4: 1,5 mol H2/mol de reagente (eletrólise 1,0; Li + H2O 0,25) & diferente & ambos \\
\bottomrule\end{longtable}%
\begin{itemize}%
\item%
\textbf{Reação 1: Li + H2O (estados, estequiometria, condições))} É a mesma reação, só com os coeficientes multiplicados por 2. Os estados físicos coincidem e as condições também (contato direto, condição ambiente). O balanço confere: Li 2=2, H 4=4, O 2=2. Conferi $\Delta$E$^\circ$ $\approx$ -0,83 - (-3,04) = +2,21 V.%
\item%
\textbf{Reação 2 (escolha livre))} O enunciado deixa explícito que vale usar um reagente sozinho ("reagindo individualmente"), e o gabarito só apresenta um exemplo ("por exemplo"). Por isso a pirólise do CH4 é tão válida quanto a eletrólise da H2O. Conferi a termodinâmica da resolução: $\Delta$S$^\circ$ = 5,74 + 2$\cdot$130,68 - 186,26 $\approx$ +80,8 J/K, então T* = $\Delta$H/$\Delta$S $\approx$ 925 K $\approx$ 652 $^\circ$C, como ela afirma. Os estados, o balanço (C 1=1, H 4=4) e as condições (sem O2, pressão baixa por Le Chatelier) estão corretos.%
\item%
\textbf{Reação 3: reforma a vapor do CH4)} A equação, os estados e a estequiometria são idênticos. Conferi o $\Delta$H$^\circ$ = -110,5 - (-74,8 - 241,8) = +206,1 kJ/mol. A única diferença nas condições é a pressão. O gabarito fala em 'alta pressão'. A resolução diz 'não muito alta' e lembra que Le Chatelier favorece pressão baixa (2 $\to$ 4 mol de gás), mas cita os 3--25 bar usados na indústria. A redação da resolução é termodinamicamente mais precisa, e as duas seriam aceitas.%
\item%
\textbf{iv) Reação que rende mais H2 por mol de reagente (reação completa))} As duas usam a mesma convenção (H2 dividido pelo total de mol de reagentes da equação balanceada) e escolhem a melhor reação dentro do próprio conjunto. Recalculei: Li + H2O 1/4 = 0,25; pirólise 2/1 = 2,0; reforma 3/2 = 1,5; eletrólise 1/1 = 1,0 (2 H2O $\to$ 2 H2 + O2). A divergência vem só da escolha da reação 2: o gabarito não incluiu a pirólise, então a reforma (1,5) é de fato a maior do conjunto dele. Se a pergunta fosse lida como 'entre todas as possíveis', a pirólise (2,0) superaria a reforma. A conclusão do gabarito vale só para o conjunto que ele escolheu, e a da resolução está certa e bem justificada para o dela. A resolução também discute outra base de comparação (por mol de CH4 a reforma rende 3; o processo global CH4 + 2H2O $\to$ CO2 + 4H2 dá 4/3 $\approx$ 1,33 por mol de reagentes).%
\end{itemize}%
\paragraph{Pontos de dúvida}%
\begin{itemize}%
\item%
\textit{(statement)} A questão é aberta: o aluno escolhe as três reações, e a resposta do item iv depende dessa escolha. Ao corrigir, o item iv deve ser avaliado pela coerência com as reações apresentadas, e não pela coincidência com a 'reforma a vapor' do gabarito. Se o corretor exigir literalmente a resposta do gabarito, a resolução (que responde pirólise) poderia ser penalizada sem razão.%
\item%
\textit{(statement)} 'Por mol de reagente' é ambíguo: pode ser o total de mol de reagentes ou o mol de uma espécie (por exemplo, o CH4). O gabarito e a resolução adotam o total. A resolução declara essa hipótese e mostra a outra leitura (por mol de CH4 a reforma vence), o que cobre os dois casos.%
\item%
\textit{(key)} Na reforma a vapor, o gabarito pede 'alta pressão'. Pelo princípio de Le Chatelier (2 $\to$ 4 mol de gás), pressão alta desfavorece a conversão; na indústria se usa pressão moderada (\textasciitilde{}20--30 bar) por razões de processo. A redação da resolução ('pressão não muito alta', 3--25 bar) é mais precisa. Isso não muda a nota.%
\item%
\textit{(key)} A eletrólise da água (reação II do gabarito) é uma reação de 'H2O reagindo individualmente', mas depende de corrente elétrica e de um eletrólito (o gabarito sugere LiOH, que não está na lista de reagentes, embora atue apenas como eletrólito). É um exemplo válido. O ponto só mostra que o gabarito também tratou o enunciado de forma flexível.%
\end{itemize}%
\section*{Questão 7}%
\href{https://poliedroresolve.sistemapoliedro.com.br/vestibulares/ita/2027/ita-2-fase-dia-2-21-10-2026-2-fase-dia-2/dissertativa-7-quimica-2-fase-dia-2-ita-2027}{Resolução do gabarito}\par\smallskip%
\begin{longtable}{@{}p{0.06\textwidth} p{0.33\textwidth} p{0.33\textwidth} p{0.1\textwidth} p{0.08\textwidth}@{}}\toprule Item & Resolução & Gabarito & Veredito & Certo \\ \midrule
a & Etapa 1: CH3CH2OH + O2 $\to$ CH3COOH + H2O (oxidação completa; oxidante em excesso, aquecimento sob refluxo, meio ácido; em laboratório K2Cr2O7/H2SO4 ou KMnO4/H+; com O2, catalisador). Também dá a equação iônica com dicromato. Etapa 2: CH3COOH + CH3CH2OH ? CH3COOCH2CH3 + H2O (esterificação de Fischer, H2SO4 conc. como catalisador, aquecimento, excesso de etanol ou retirada de água/éster) & 1ª etapa: C2H5OH + O2 $\to$ CH3COOH + H2O (K2Cr2O7 acidificado com H2SO4, a quente, refluxo); 2ª etapa: CH3COOH + C2H5OH ? CH3COOC2H5 + H2O (H2SO4, aquecimento, refluxo) & igual & ambos \\
b & $\Delta$H$^\circ$ = -496,2 kJ/mol & $\Delta$H$^\circ$ = -496,2 kJ/mol & igual & ambos \\
c & $\Delta$S$^\circ$ = -127,2 J/(mol$\cdot$K) & $\Delta$S$^\circ$ = -127,2 J/(K$\cdot$mol) & igual & ambos \\
d & $\Delta$G$^\circ$ $\approx$ -458,3 kJ/mol (espontânea a 298 K) & $\Delta$G$^\circ$ $\approx$ -458,3 kJ/mol & igual & ambos \\
\bottomrule\end{longtable}%
\begin{itemize}%
\item%
\textbf{a)} As equações e as condições são as mesmas do gabarito: oxidante forte acidificado, a quente e sob refluxo na 1ª etapa; catálise por H2SO4 e aquecimento na 2ª. A soma das etapas dá a reação global. A equação iônica extra está correta em átomos (C 6, H 34, O 17, Cr 4) e em carga (+12 nos dois lados). O balanço de elétrons também confere: 4 e$^{-}$ por etanol e 6 e$^{-}$ por Cr2O7$^2$$^{-}$.%
\item%
\textbf{b)} Recalculado: {[}-480,0 + 2(-285,8){]} - {[}2(-277,7) + 0{]} = -1051,6 + 555,4 = -496,2 kJ/mol.%
\item%
\textbf{c)} Recalculado: {[}259,4 + 2(69,9){]} - {[}2(160,7) + 205,0{]} = 399,2 - 526,4 = -127,2 J$\cdot$mol$^{-}$$^{1}$$\cdot$K$^{-}$$^{1}$.%
\item%
\textbf{d)} $\Delta$G$^\circ$ = -496,2 - 298(-0,1272) = -496,2 + 37,9 = -458,29 $\approx$ -458,3 kJ/mol. A resolução converteu $\Delta$S$^\circ$ para kJ corretamente. A conclusão sobre a espontaneidade é um acréscimo coerente.%
\end{itemize}%
\paragraph{Pontos de dúvida}%
\begin{itemize}%
\item%
\textit{(solution)} No item a, a resolução escreve a etapa 1 com O2 (para que a soma das etapas dê a reação global), mas as condições citadas misturam dois cenários: dicromato/H+ em laboratório e O2 com catalisador ou Acetobacter. Está correto e é o mesmo que o gabarito faz (equação com O2, condição com K2Cr2O7/H2SO4). Um corretor rigoroso poderia, porém, pedir coerência entre o oxidante da equação e o das condições.%
\item%
\textit{(key)} Na resolução do item d, o gabarito escreve a unidade final de $\Delta$G como 'kJ', e não 'kJ/mol'. É só notação; o valor -458,3 está correto.%
\end{itemize}%
\section*{Questão 8}%
\href{https://poliedroresolve.sistemapoliedro.com.br/vestibulares/ita/2027/ita-2-fase-dia-2-21-10-2026-2-fase-dia-2/dissertativa-8-quimica-2-fase-dia-2-ita-2027}{Resolução do gabarito}\par\smallskip%
\begin{longtable}{@{}p{0.06\textwidth} p{0.33\textwidth} p{0.33\textwidth} p{0.1\textwidth} p{0.08\textwidth}@{}}\toprule Item & Resolução & Gabarito & Veredito & Certo \\ \midrule
a & S (Z=16): 2 elétrons desemparelhados; Se (Z=34): 2 elétrons desemparelhados; K (Z=19): 1 elétron de valência (4s$^{1}$) & S: 2 desemparelhados; Se: 2 desemparelhados; K: 1 elétron na camada de valência & igual & ambos \\
b & H2S: 2 pares não ligantes no S; SeF4: 1 par não ligante no Se & H2S: 2 pares não ligantes no S; SeF4: 1 par não ligante no Se & igual & ambos \\
c & H2S: angular (em V); SeF4: gangorra (seesaw) & H2S: angular; SeF4: gangorra & igual & ambos \\
\bottomrule\end{longtable}%
\begin{itemize}%
\item%
\textbf{a)} S: {[}Ne{]}3s$^2$3p$^{4}$ e Se: {[}Ar{]}3d$^{1}$$^{0}$4s$^2$4p$^{4}$. Pela regra de Hund, o p$^{4}$ fica ?? ? ?, então são 2 elétrons desemparelhados em cada. K: {[}Ar{]}4s$^{1}$, com 1 elétron na camada de valência. A resolução e o gabarito dão o mesmo resultado, e está certo.%
\item%
\textbf{b)} S e Se têm 6 elétrons de valência. No H2S o S faz 2 ligações e sobram (6-2)/2 = 2 pares livres (AX2E2). No SeF4 o Se faz 4 ligações e sobra (6-4)/2 = 1 par livre (AX4E). Os compostos foram bem identificados (o enunciado pede um hidreto binário estável do S e um composto do Se com quatro ligações simples a F) e a contagem bate com o gabarito.%
\item%
\textbf{c)} No AX2E2, o arranjo dos pares é tetraédrico e a geometria molecular é angular. No AX4E, o arranjo é bipirâmide trigonal com o par livre na posição equatorial, o que dá a geometria em gangorra. Coincide com o gabarito. As justificativas extras da resolução também estão corretas: ângulo H--S--H menor que 109,5$^\circ$ e ângulos F--Se--F menores que 180$^\circ$ e 120$^\circ$.%
\end{itemize}%
\paragraph{Pontos de dúvida}%
\begin{itemize}%
\item%
\textit{(solution)} Na conferência, a resolução diz que a expansão do octeto é "permitida para elementos do 4º período". A afirmação é verdadeira para o Se, mas a explicação usual vale do 3º período em diante. Isso não afeta nenhuma resposta.%
\item%
\textit{(key)} A resolução do gabarito no item a) usa imagens com nomes de outra questão (Q03\_A2, Q03\_A3), o que parece reaproveitamento ou troca de figuras na página do Poliedro. Não deu para ver as imagens, mas as respostas por item do gabarito estão completas e corretas.%
\end{itemize}%
\section*{Questão 9}%
\href{https://poliedroresolve.sistemapoliedro.com.br/vestibulares/ita/2027/ita-2-fase-dia-2-21-10-2026-2-fase-dia-2/dissertativa-9-quimica-2-fase-dia-2-ita-2027}{Resolução do gabarito}\par\smallskip%
\begin{longtable}{@{}p{0.06\textwidth} p{0.33\textwidth} p{0.33\textwidth} p{0.1\textwidth} p{0.08\textwidth}@{}}\toprule Item & Resolução & Gabarito & Veredito & Certo \\ \midrule
a & {[}HX{]} = 8$\cdot$10$^{-}$$^3$ mol$\cdot$L$^{-}$$^{1}$ (0,008) & {[}HX{]} = 8,0$\cdot$10$^{-}$$^3$ mol$\cdot$L$^{-}$$^{1}$ & igual & ambos \\
b & {[}HY{]} = 1,2$\cdot$10$^{-}$$^2$ mol$\cdot$L$^{-}$$^{1}$ (0,012) & {[}HY{]} = 1,2$\cdot$10$^{-}$$^2$ mol$\cdot$L$^{-}$$^{1}$ & igual & ambos \\
\bottomrule\end{longtable}%
\begin{itemize}%
\item%
\textbf{a)} A base gasta é n(ZOH) = 0,1$\cdot$0,020 = 2,0$\cdot$10$^{-}$$^3$ mol, que neutraliza HX e HY juntos. Tirando n(HY) = 1,2$\cdot$10$^{-}$$^3$ mol, sobra n(HX) = 8$\cdot$10$^{-}$$^{4}$ mol. Em 0,100 L, isso dá 8$\cdot$10$^{-}$$^3$ mol/L. A resolução segue o mesmo raciocínio do gabarito e chega ao mesmo valor.%
\item%
\textbf{b)} No ponto de equivalência, pH = 9 dá {[}OH$^{-}${]} = 10$^{-}$$^{5}$. Com Kh = Kw/Ka = 10$^{-}$$^{8}$, sai {[}Y$^{-}${]} = 10$^{-}$$^{1}$$^{0}$/10$^{-}$$^{8}$ = 10$^{-}$$^2$ mol/L em 120 mL. Voltando ao volume inicial: c\_Y = 10$^{-}$$^2$$\cdot$120/100 = 1,2$\cdot$10$^{-}$$^2$ mol/L. Também resolvi sem aproximações a partir de pH = 9 e obtive c\_Y = 0,012011 e c\_X = 0,007989. No sentido inverso, partindo de c\_Y = 0,012, o pH no equivalente sai 8,9998. Isso confirma as duas respostas.%
\end{itemize}%
\paragraph{Pontos de dúvida}%
\begin{itemize}%
\item%
\textit{(key)} Erro de digitação na resolução do gabarito: escreve 'no equilíbrio, {[}OH$^{-}${]} = {[}HX{]}', quando o certo é {[}OH$^{-}${]} = {[}HY{]} (produtos da hidrólise de Y$^{-}$). A conta seguinte usa o valor certo (10$^{-}$$^{5}$)$^2$, então o resultado não muda.%
\item%
\textit{(solution)} Na conferência, a resolução chama de 'balanço de cargas exato' a equação {[}H?{]} + {[}HY{]} = {[}OH$^{-}${]}. Essa equação é o balanço de prótons (ou a combinação do balanço de cargas com o de massa), como a própria resolução dizia no Passo 2. É só uma questão de nome, sem efeito no resultado.%
\item%
\textit{(statement)} O enunciado não informa a temperatura nem Kw. A resolução e o gabarito supõem 25 $^\circ$C (Kw = 10$^{-}$$^{1}$$^{4}$), o que é razoável e foi declarado nos dois. O enunciado também não diz qual ponto de equivalência é. A resolução justifica bem por que tem de ser o equivalente total (HX + HY): com pH 9 não pode sobrar HY sem neutralizar.%
\end{itemize}%
\section*{Questão 10}%
\href{https://poliedroresolve.sistemapoliedro.com.br/vestibulares/ita/2027/ita-2-fase-dia-2-21-10-2026-2-fase-dia-2/dissertativa-10-quimica-2-fase-dia-2-ita-2027}{Resolução do gabarito}\par\smallskip%
\begin{longtable}{@{}p{0.06\textwidth} p{0.33\textwidth} p{0.33\textwidth} p{0.1\textwidth} p{0.08\textwidth}@{}}\toprule Item & Resolução & Gabarito & Veredito & Certo \\ \midrule
a(1) & X ({[}Fe2+{]} alta) é o cátodo e Y é o ânodo (pilha de concentração iônica, E = E$^\circ$ + 0,0295 log{[}Fe2+{]}) & X catódica, Y anódica & igual & ambos \\
a(2) & Y (pH baixo) é o cátodo e X (pH alto) é o ânodo, pelo modelo de Nernst com Fe(OH)2 (E\_Fe $\approx$ -0,1 - 0,059 pH). A própria resolução admite que a passivação inverteria isso. & X (pH alto, OH-) catódica e Y (pH baixo, H+) anódica & diferente & gabarito \\
a(3) & X ({[}O2{]} alta) é o cátodo e Y ({[}O2{]} baixa) é o ânodo (aeração diferencial) & X catódica, Y anódica (aeração diferencial) & igual & ambos \\
a(4) & X ({[}Cl-{]} alta) é o ânodo e Y é o cátodo (despassivação e pite autocatalítico; complexação de Fe2+ como efeito menor) & X (Cl- alto) anódica, Y catódica (pite por cloreto) & igual & ambos \\
b & Proteção catódica nos quatro casos (ânodo de sacrifício como Zn/Al/Mg, ou corrente impressa). A anódica é inadequada porque o filme passivo do Fe é instável em meio neutro ou ácido e é destruído pelo Cl-. & Proteção catódica em todos os casos (1 a 4). A película passiva do ferro é frágil, então a proteção anódica não é adequada. & igual & ambos \\
\bottomrule\end{longtable}%
\begin{itemize}%
\item%
\textbf{a(1))} Pilha de concentração iônica diferencial: o potencial de redução de Fe2+/Fe é maior onde {[}Fe2+{]} é maior, então ali é o cátodo. As duas atribuições são iguais. A conta de verificação também está certa: razão 100 dá 0,059 V.%
\item%
\textbf{a(2))} A atribuição sai invertida. As contas da resolução estão certas: com Ksp = 5$\cdot$10\^{}{-}17, E\_Fe = -0,44 + 0,0295$\cdot$log(Ksp/{[}OH-{]}$^2$) $\approx$ -0,10 - 0,059$\cdot$pH, e por exemplo pH 12 dá -0,80 V. O que falha é a interpretação. Esse potencial mais baixo em pH alto vem da formação de um produto sólido (Fe(OH)2 e depois Fe3O4/Fe2O3). Pelo diagrama de Pourbaix, é a região de passivação, e não de dissolução ativa. Além disso, os potenciais do ferro e do oxidante caem os dois 0,059 V por unidade de pH, então a força motriz em cada região é a mesma. Quem decide é a cinética e a passivação. Em meio ácido o Fe se dissolve ativamente, com H+ reduzindo junto. Em meio alcalino ele fica passivo e com potencial mais nobre. É o mesmo princípio da pilha clássica aço{-}no{-}concreto (cátodo) contra aço{-}no{-}solo (ânodo), que a própria resolução cita na ressalva. A leitura que o enunciado espera (e que o gabarito dá) é que a região ácida Y corrói, ou seja, Y é o ânodo e X o cátodo. A resolução mencionou essa possibilidade, mas escolheu a resposta contrária.%
\item%
\textbf{a(3))} Aeração diferencial: a região mais oxigenada é o cátodo (O2 + 2H2O + 4e- $\to$ 4OH-) e a perda de metal fica na região menos aerada. A conta de 0,0295 V para razão 100 em pO2 está certa.%
\item%
\textbf{a(4))} Mesma atribuição e mesmo mecanismo: o Cl- rompe o filme passivo, o Fe2+ hidrolisa e acidifica o pite, que atrai mais Cl-. O argumento extra de que a formação de FeCl+ reduz o Fe2+ livre é correto, mas secundário.%
\item%
\textbf{b)} Mesma conclusão e justificativa compatível. Os potenciais citados estão certos: Fe -0,44 V, Zn -0,76 V, Al -1,66 V, Mg -2,37 V.%
\end{itemize}%
\paragraph{Pontos de dúvida}%
\begin{itemize}%
\item%
\textit{(solution)} No item a(2) a resolução trocou cátodo e ânodo. Ela usou só o potencial de equilíbrio de Nernst com Fe(OH)2 e deixou de fora a passivação do Fe em meio alcalino e a dissolução ativa em meio ácido. Na leitura padrão de corrosão, a região de pH baixo (Y) é o ânodo, como dá o gabarito. A resolução cita essa inversão na ressalva, mas fica com a resposta contrária, o que deve custar a nota do subitem.%
\item%
\textit{(statement)} No caso (2) o enunciado não diz se o meio está aerado nem qual é a faixa de pH. Num modelo puramente termodinâmico, sem O2 e sem passivação, dá para defender X como ânodo (o potencial de corrosão em meio alcalino desaerado é mais negativo). Por isso a questão admite discussão, mas a resposta esperada no vestibular é a do gabarito.%
\item%
\textit{(key)} No item b, a frase "para cada caso (itens 1 a 4)" sugere que a banca talvez aceitasse ou esperasse respostas diferentes por caso. Por exemplo, proteção anódica na região alcalina de (2), já que o aço{-}carbono pode ser protegido anodicamente em meios alcalinos e em H2SO4 concentrado. O gabarito do Poliedro e a resolução concordam em proteção catódica para todos os casos. Vale conferir quando sair o gabarito oficial.%
\item%
\textit{(key)} A justificativa do gabarito em a(2) ("OH- desfavorece a redução de H+, logo o ferro oxida menos") é cinética e resumida. Não menciona a passivação em pH alto nem diz qual reação catódica ocorre em X (redução de O2). A conclusão está certa, mas a explicação está incompleta.%
\end{itemize}%
\end{document}